\section{Opis wyników prac}

Na obecnym etapie wykonano:

\begin{itemize}
	\item sformułowanie tematu i zakresu projektu,
	\item przegląd podobnych rozwiązań i stanu wiedzy, w tym ewolucji metod detekcji,
	\item uzupełnienie tła o komercyjne znakowanie wodne OpenAI (textGrain) i kontekst \textit{EU AI Act},
	\item konsultacje z ekspertem (dr hab.\ Piotr Hońko) oraz z narzędziem GenAI,
	\item porównanie obu źródeł rekomendacji,
	\item uszczegółowienie wymagań funkcjonalnych i niefunkcjonalnych,
	\item doprecyzowanie zakresu, kamieni milowych i harmonogramu (diagram Gantta),
	\item szczegółową specyfikację wejścia/wyjścia, przepływu przetwarzania i planu ewaluacji,
	\item wstępną listę zbiorów danych (Kaggle, SpeakLeash i inne).
\end{itemize}

Implementacja metod detekcji oraz pomiary dokładności będą przedmiotem kolejnych etapów projektu.

\subsection{Ścieżka machine learning}

Zaimplementowano pipeline fine-tuningu modelu \path{distilbert-base-uncased}
(67 mln parametrów) w~PyTorch Lightning. Klasyfikacja jest binarna:
\texttt{human} = 0, \texttt{ai} = 1. Eksperyment używa jednego seeda (\texttt{42})
i~jednego podziału danych.

\paragraph{Czyszczenie danych i przeciek.}
Przed podziałem usunięto dokładne duplikaty oraz teksty krótsze niż 5 słów.
Z 29\,145 wierszy pozostało 27\,296. Podział train/validation/test wynosi
21\,838 / 2\,729 / 2\,729, a liczba wspólnych tekstów między splitami wynosi
\textbf{0}. Wcześniejsze wyniki były zawyżone przez duplikaty: przed poprawką
10,4\% testu miało dokładny odpowiednik w zbiorze treningowym.

\paragraph{Strojenie hiperparametrów.}
Uruchomiono jawną siatkę czterech konfiguracji. Zbiór testowy nie był używany
podczas strojenia; wyboru dokonano na podstawie \texttt{val\_f1}.

\begin{center}
\begin{tabular}{|c|c|c|c|c|}
\hline
learning rate & batch size & max\_length & best val F1 & czas \\
\hline
2e-5 & 32 & 256 & \textbf{0,99821} & 186 s \\
3e-5 & 32 & 256 & 0,99821 & 195 s \\
4e-5 & 64 & 256 & 0,99821 & 212 s \\
2e-5 & 32 & 128 & 0,99599 & 157 s \\
\hline
\end{tabular}
\end{center}

Wybrano konfigurację \textbf{lr = 2e-5, batch size = 32, max\_length = 256}.
Liczba epok jest kontrolowana przez early stopping: próby mają limit 4 epok,
a finalny trening — 20 epok z \texttt{patience = 5} na \texttt{val\_f1}.

\paragraph{Wynik finalny.}
Po 12 epokach (najlepszy checkpoint z epoki 7) uzyskano na odłożonym zbiorze
testowym:

\begin{center}
\begin{tabular}{|l|c|}
\hline
metryka & wartość \\
\hline
accuracy & \textbf{0,99707} \\
precision & \textbf{1,00000} \\
recall & \textbf{0,99284} \\
F1 & \textbf{0,99641} \\
\hline
\end{tabular}
\end{center}

Macierz pomyłek: \texttt{[[1611, 0], [8, 1110]]}, czyli 8 błędów na 2\,729
próbek, wszystkie w stronę \texttt{ai} \(\rightarrow\) \texttt{human}.
Cały finalny trening zajął 8,4 min.

\paragraph{Stabilność i koszt.}
Krzywa uczenia jest stabilna: najgorszy epizod to epoka 3
(\texttt{val\_f1} = 0,9933, \texttt{val\_loss} = 0,0172). Fine-tuning wymaga
GPU; na karcie RX 9070 XT (ROCm, bf16) jedna epoka zajmuje ok. 40 s, a pełny
trening 8,4 min. Starszy wariant BERT-base we fp32 potrzebował ok. 9 min na
epokę.

\paragraph{Zależność od długości tekstu.}
Średnia długość tekstu human to ok. 430 słów, a AI — ok. 312 słów. Sam
klasyfikator oparty wyłącznie na liczbie słów osiąga AUC = 0,691, co oznacza,
że długość jest istotnym confoundem. Dlatego metryki uzupełniono rozbiciem na
kwartyle długości:

\begin{center}
\begin{tabular}{|l|c|c|c|c|c|}
\hline
kwartyl (słowa) & próbki & accuracy & precision & recall & F1 \\
\hline
Q1 (61--251) & 692 & 0,99855 & 1,00000 & 0,99746 & 0,99873 \\
Q2 (251--345) & 678 & 0,99410 & 1,00000 & 0,98892 & 0,99443 \\
Q3 (345--464) & 677 & 0,99557 & 1,00000 & 0,98790 & 0,99391 \\
Q4 (464--1436) & 682 & 1,00000 & 1,00000 & 1,00000 & 1,00000 \\
\hline
\end{tabular}
\end{center}

\paragraph{Ograniczenia metody ML.}
Do najważniejszych ograniczeń należą: koszt obliczeniowy fine-tuningu (wymagany GPU),
overfitting (\texttt{train\_loss} ok. \(1{,}3 \cdot 10^{-6}\) przy
\texttt{val\_loss} ok. 0,006), zależność od domeny (angielskie eseje
z~jednego zbioru, brak zewnętrznego testu), zależność od długości tekstu,
pozostałe near-duplikaty oraz nierównowaga klas. Model nie był kalibrowany,
a~etykiety nie były weryfikowane ręcznie. Pełny opis znajduje się w~dokumencie \path{train/ml-ograniczenia.md}.

\subsection{Ścieżka statystyczna}

Zaimplementowano interpretowalny pipeline statystyczno-stylometryczny bazujący
na wektorze 13 cech leksykalno-syntaktycznych. Ewaluację przeprowadzono na
dokładnie tym samym oczyszczonym zbiorze (27\,280 tekstów) i identycznym
podziale danych (\texttt{seed = 42}, 21\,824 / 2\,728 / 2\,728) co w ścieżce ML,
co zapewnia pełną spójność metodologiczną.

\paragraph{Ekstrakcja cech stylometrycznych.}
Dla każdego tekstu obliczono 13 cech odzwierciedlających rytm i bogactwo językowe:
\begin{itemize}
	\item \textbf{Długość zdań} (\path{avg_sentence_len}, \path{std_sentence_len}) --
	średnia liczba słów w~zdaniu oraz jej odchylenie standardowe. Teksty ludzkie charakteryzują się średnio dłuższymi zdaniami (25,9 vs 23,0 słów) oraz znacznie większą zmiennością (odchylenie 13,0 vs 8,4).
	\item \textbf{Burstiness (zmienność rytmu)} --
	cecha mierzona współczynnikiem zmienności $CV = \sigma / \mu$ (\path{burstiness_cv}: 0,491 dla ludzi vs 0,353 dla AI) oraz średnią bezwzględną różnicą długości kolejnych zdań (\path{burstiness_diff}: 15,33 vs 9,00). Tekst ludzki cechuje wyraźna dynamikę skokowa, podczas gdy modele LLM tworzą zdania o~wygładzonym rytmie.
	\item \textbf{Bogactwo słownikowe} --
	wskaźnik Type-Token Ratio $\text{TTR} = |V| / N$ (\path{ttr}) oraz indeks Guirauda (\path{root_ttr} = $|V| / \sqrt{N}$).
	\item \textbf{Cechy leksykalne i interpunkcyjne} --
	udział słów funkcyjnych (\path{stopword_ratio}), średnia długość słowa (\path{avg_word_len}: 4,42 dla ludzi vs 5,11 dla AI), gęstość przecinków na zdanie (\path{comma_per_sentence}: 0,83 vs 0,99) oraz proporcja wielkich liter i~cyfr.
\end{itemize}

\paragraph{Modele klasyfikacyjne.}
Zbadano dwa warianty decyzyjne operujące na wyekstrahowanych cechach:
\begin{enumerate}
	\item \textbf{Model progowy (drzewo o głębokości 2):} jawne reguły decyzyjne oparte o średnią długość wyrazu i burstiness osiągnęły dokładność \textbf{85,34\%} oraz F1 = \textbf{0,8296}.
	\item \textbf{Regresja logistyczna ze standaryzacją (\texttt{StandardScaler}):} model liniowy osiągnął wynik finalny na odłożonym zbiorze testowym:
\end{enumerate}

\begin{center}
\begin{tabular}{|l|c|}
\hline
metryka & wartość \\
\hline
accuracy & \textbf{0,90689} \\
precision & \textbf{0,86986} \\
recall & \textbf{0,90877} \\
F1 & \textbf{0,88889} \\
ROC-AUC & \textbf{0,96836} \\
\hline
\end{tabular}
\end{center}

Macierz pomyłek: \texttt{[[1458, 152], [102, 1016]]}. Na 2\,728 próbek testowych
model popełnił 254 błędy (152 fałszywe alarmy $\text{human} \rightarrow \text{ai}$ oraz
102 błędy $\text{ai} \rightarrow \text{human}$).

\paragraph{Interpretowalność wag cech.}
Analiza znormalizowanych współczynników regresji logistycznej wykazała, że
najsilniejszymi predyktorami tekstu ludzkiego są: średnia długość zdań (waga $-2{,}74$),
skokowość długości sąsiednich zdań \path{burstiness_diff} (waga $-2{,}56$),
zmienność rytmu \path{burstiness_cv} (waga $-1{,}72$) oraz udział wielkich liter (waga $-1{,}45$).
Z kolei za tekstem AI najsilniej przemawiają: średnia długość wyrazu (waga $+3{,}08$,
sugerująca używanie bardziej wyszukanego leksykonu), odchylenie długości zdań (waga $+3{,}05$)
oraz zagęszczenie przecinków (waga $+0{,}94$).

\paragraph{Zależność od długości tekstu.}
Podział odłożonego zbioru testowego na kwartyle długości słownej ujawnia wyraźną
wrażliwość metody statystycznej na objętość materiału:

\begin{center}
\begin{tabular}{|l|c|c|c|c|c|}
\hline
kwartyl (słowa) & próbki & accuracy & precision & recall & F1 \\
\hline
Q1 (50--256) & 685 & 0,84088 & 0,83846 & 0,87668 & 0,85714 \\
Q2 (256--345) & 680 & 0,89559 & 0,90000 & 0,90736 & 0,90366 \\
Q3 (345--453) & 681 & 0,93098 & 0,87755 & 0,95911 & 0,91652 \\
Q4 (453--1650) & 682 & \textbf{0,96041} & 0,85965 & 0,89908 & 0,87892 \\
\hline
\end{tabular}
\end{center}

Na krótkich tekstach (Q1) dokładność wynosi 84,1\%, podczas gdy na długich esejach
(Q4) wzrasta do 96,0\%. Wynika to z dużej wariancji estymatorów statystycznych
(TTR, burstiness) przy małej liczbie zdań.

\paragraph{Ograniczenia metody statystycznej.}
Kluczowe ograniczenia to: wrażliwość na małą długość tekstu (poniżej 150 słów
statystyki ulegają degradacji), podatność na celowy prompt engineering
(instrukcja nakazująca zróżnicowanie zdań znosi sygnaturę stylu), brak analizy
semantyki i logiki oraz konieczność rekalibracji przy zmianie języka (np. na język polski).
Pełną analizę zawarto w dokumencie \path{statistical/stat-ograniczenia.md}.

\subsection{Porównanie ścieżki statystycznej i machine learning}

Dzięki zachowaniu identycznego podziału danych (\texttt{seed = 42}) możliwe
było bezpośrednie zestawienie obu podejść na odłożonym zbiorze testowym:

\begin{center}
\begin{tabular}{|l|c|c|}
\hline
Cecha / Metryka & Ścieżka statystyczna & Ścieżka ML (DistilBERT) \\
\hline
accuracy & 0,90689 & \textbf{0,99707} \\
precision & 0,86986 & \textbf{1,00000} \\
recall & 0,90877 & \textbf{0,99284} \\
F1-score & 0,88889 & \textbf{0,99641} \\
liczba błędów (na 2\,728) & 254 & \textbf{8} \\
czas trenowania & \textbf{ok. 25 s (CPU)} & ok. 8,4 min (GPU) \\
liczba parametrów & \textbf{14} & 67 mln \\
interpretowalność & \textbf{pełna (jawne wagi)} & niska (czarna skrzynka) \\
\hline
\end{tabular}
\end{center}

Model transformerowy osiąga niemal bezbłędną separację (F1 $\approx 0{,}996$,
brak fałszywych alarmów), jednak wymaga kosztownego sprzętu obliczeniowego.
Ścieżka statystyczna osiąga satysfakcjonującą jakość (dokładność $> 90\%$, F1 $\approx 0{,}89$,
znacznie powyżej założonego w wymaganiach minimum 51\%), oferując natychmiastowy
czas wnioskowania oraz pełną przejrzystość podjętej decyzji.
